\chapter{Evaluation Methodology}\label{ch:methodology}

This chapter describes how measurements are taken; the concrete results follow in
Chapter~\ref{ch:results}.

\section{Three Measurement Series}\label{sec:series}

The architecture prescribes three measurement series, all executed by the
\texttt{messung\_driver} through the \texttt{ExperimentDriver} library, configured either
as three fixed series or via an external \texttt{messreihen.xml}.

\begin{description}
  \item[Series A --- PRT-ART vs.\ the state of the art.] Compares the probe against the
        documented SOTA algorithms. Two modes: \emph{A\_defined} (quick verification
        against the 8 Tier-1 profiles) and \emph{A\_full} (all 30 SOTA profiles; the
        default, per the directive to measure as completely as possible).
  \item[Series B --- cache-engine permutations.] The pure ``state of the art'' baseline:
        all SOTA profiles without the PRT-ART probe.
  \item[Series C --- merge of old and new.] Concrete merge points between PRT-ART building
        blocks and cache-engine permutations, e.g.\ PRT-ART OLC + tcmalloc + ART page;
        PRT-ART 4+2 pool + HOT page; PRT-ART path prefetch + B\textsuperscript{2}-tree
        page; PRT-ART multi-level layout + Wormhole page.
\end{description}

\section{ExperimentDriver Phases}\label{sec:experimentdriver}

The \texttt{ExperimentDriver} runs seven phases per series: (1)~\textbf{enumerate} (parse
the XML configs and build the permutation list); (2)~\textbf{codegen} (one
\texttt{comdare\_perm\_<fp>.cpp} per permutation, plus
\texttt{generate\_module\_from\_profile} for the SOTA profiles); (3)~\textbf{compile};
(4)~\textbf{load} (LoadLibrary/dlopen of all modules); (5)~\textbf{execute}
(\texttt{create\_instance} + profile-aware \texttt{run\_workload}); (6)~\textbf{measure}
(aggregate the binary measurement records); (7)~\textbf{persist}
(\texttt{measurements.csv} + \texttt{.json}). Two opt-in phases hot-compile missing modules
and verify the ABI contract per module.

\section{Hypotheses}\label{sec:hypotheses}

\begin{description}
  \item[H1 (page-type cost)] The mean access cost per page-type variant differs
        systematically by workload: dense page types are expected to win under Zipfian
        (YCSB-A) and range-optimized pages under range scans (YCSB-E).
  \item[H2 (code quality per source)] The code-quality score inherited from the original
        source repository correlates measurably with achieved throughput.
  \item[H3 (inline vs.\ external vs.\ chain distribution)] The observed distribution of the
        three value-handle variants correlates with page density: dense pages favour
        inline, sparse pages external, and PRT-ART heuristics chain-ref.
\end{description}

\section{Workloads and Datasets}\label{sec:workloads}

Phase~5 routes the YCSB workload per permutation from the profile's traversal tag
(Table~\ref{tab:workload-routing}); modules without a profile use the global default
workload.

\begin{table}[h]
\centering
\caption{Profile-aware workload routing}\label{tab:workload-routing}
\begin{tabular}{lll}
\toprule
\textbf{Traversal tag} & \textbf{YCSB workload} & \textbf{Characteristic} \\
\midrule
\texttt{RANGE\_SCAN} / \texttt{RANGE\_HOT\_PATH} & YCSB-E & range queries \\
\texttt{HOT\_PATH} / \texttt{PRTART\_HOT\_PATH} & YCSB-A & 50/50 read/update, Zipfian \\
\texttt{STANDARD} (default) & YCSB-C & read-only \\
\bottomrule
\end{tabular}
\end{table}

\section{Experimental Platforms}\label{sec:platforms}

Measurements target a local workstation for development verification and the TU~Dresden
ZIH cluster (Barnard CPU, Capella GPU) for the full sweeps, plus the Tier-3 platforms
(Ryzen~9~9950X3D, Intel~i9~14900KS, and an HBM platform). An \texttt{IPlatformProbe}
(phase~1) reports, per platform, the runnable subset of axes (e.g.\ AVX-512 only where
available, HBM-aware only on HBM hardware).

\section{Permutation Explosion and Reduction}\label{sec:explosion}

With the N-phase extension from 11 to 14 axes plus $\sim$30 sub-axes, the space grows from
$\sim$5.5 billion to more than $10^{11}$ permutations. Three mechanisms tame the
explosion: (1)~the \texttt{<expected\_workload>} profile filter routes only compatible
workloads; (2)~\texttt{MessreihenMode::Defined} runs only the declared tuples
($\sim$100--1000 per series) that feed the manuscript; (3)~\texttt{Full} (and the more
realistic \texttt{Full-Sampled}, e.g.\ every 1000th permutation) runs only on the ZIH
cluster, respecting axis compatibility reported by the platform probe.
